\section*{अध्याय \ref{ch:b3:water-rock-interaction} --- जल--चट्टान अंतर्क्रिया}

\begin{solution}{exo:b3:water-rock-interaction:1}
आयन: $2\times48/40.08 + 2\times9.7/24.31 + 23/22.99 + 3.9/39.10 =
\qty{4.29}{meq/L}$; ऋणायन: $183/61.02 + 35.5/35.45 + 2\times48/96.06 =
\qty{5.00}{meq/L}$। त्रुटि $100\times(4.29 - 5.00)/9.29 = \qty{-7.6}{\percent}$:
ऐसे जल के सावधानीपूर्ण विश्लेषण की \qty{1}{\percent} या \qty{2}{\percent} की सामान्य सीमा से कहीं अधिक।
कोई धनायन छूट गया है या कम दर्ज हुआ है, या कोई ऋणायन अधिक है।
\end{solution}

\begin{solution}{exo:b3:water-rock-interaction:2}
$\sqrt I = 0.0707$। \ce{Ca^2+}: $-\log\gamma = 0.5085\times4\times0.0707/(1 +
0.3281\times6\times0.0707) = 0.126$, $\gamma = 0.748$; \ce{HCO3-}: $\gamma =
0.927$; \ce{SO4^2-}: $\gamma = 0.739$।
\end{solution}

\begin{solution}{exo:b3:water-rock-interaction:3}
$\sqrt I = 0.126$; $\gamma_{\ce{Ca}} = 0.622$, $\gamma_{\ce{SO4}} = 0.602$;
$\log\mathrm{IAP} = \log(3.0\times10^{-3}\times0.622\times4.0\times10^{-3}
\times0.602) = -5.35$; $\mathrm{SI} = -5.35 + 4.55 = -0.80$। जल अल्पसंतृप्त है
और अब भी जिप्सम को घोल सकता है।
\end{solution}

\begin{solution}{exo:b3:water-rock-interaction:4}
$m = 0.012/60.08 = \qty{2.00e-4}{mol/L}$, $\log m = -3.70$। क्वार्ट्ज़:
$\mathrm{SI} = -3.70 + 3.98 = +0.28$, कुछ अति-संतृप्त; अनाकार
सिलिका: $-3.70 + 2.71 = -0.99$, अल्पसंतृप्त।
\end{solution}

\begin{solution}{exo:b3:water-rock-interaction:5}
\qty{25}{\celsius} पर, $\log K = -8.45 - 6.35 - 1.47 + 10.33 = -5.94$, $m =
(10^{-5.94}\times10^{-2}/4)^{1/3} = \qty{1.42}{mmol/L}$, \qty{57}{mg/L}; \qty{10}{\celsius}
पर, $\log K = -5.59$, $m = \qty{1.86}{mmol/L}$,
\qty{75}{mg/L}: ठंडे में एक-तिहाई अधिक।
\end{solution}

\begin{solution}{exo:b3:water-rock-interaction:6}
$[\ce{CO2}] = 10^{-7.60}\times4.0\times10^{-3}/10^{-6.35} =
\qty{2.25e-4}{mol/L}$; $[\ce{CO3^2-}] = 10^{-10.33}\times4.0\times10^{-3}/
10^{-7.60} = \qty{7.4e-6}{mol/L}$; $\log P_{\ce{CO2}} = -3.648 + 1.47 =
-2.18$, वायु से 16 गुना। कुल \qty{4.23}{mmol/L}: \ce{CO2}
\qty{5.3}{\percent}, \ce{HCO3-} \qty{94.5}{\percent}, \ce{CO3^2-}
\qty{0.18}{\percent}।
\end{solution}

\begin{solution}{exo:b3:water-rock-interaction:7}
$r = 10^{-12.09} = \qty{8.2e-13}{mol.m^{-2}.s^{-1}}$ (अम्लीय, उदासीन और
क्षारीय पद क्रमशः $10^{-12.45}$, $10^{-12.56}$ और $10^{-12.74}$ का योगदान देते हैं)। त्रिज्या
$rV_m = 8.2\times10^{-13}\times1.0131\times10^{-4} =
\qty{8.3e-17}{m/s}$ पर घटती है, यानी प्रति वर्ष \qty{2.6}{nm}; \qty{0.5}{mm} त्रिज्या
$6.0\times10^{12}$~s, लगभग \num{190000} वर्ष तक चलती है।
\end{solution}

\begin{solution}{exo:b3:water-rock-interaction:8}
$\exp[(69.8/0.008314)(1/278.15 - 1/298.15)] = 7.6$; \qty{74}{kJ/mol}
के साथ, 8.6।
\end{solution}

\begin{solution}{exo:b3:water-rock-interaction:9}
$(12.703 - X - 3Y)/3$ और $(7.435 - 2Y)/2$ को बराबर करने पर $Y$ हट जाता है: $X =
12.703 - 11.153 = 1.55$, एक क्षैतिज रेखा। जल में $X = \log(2.0\times
10^{-3}/39.10) + 7.0 = 2.71$ और $Y = \log(0.015/60.08) = -3.60$ है: 1.55 से ऊपर
और मस्कोवाइट--K-फ़ेल्सपार रेखा ($X = -3.161 - 3Y = 7.65$) से बहुत नीचे,
इसलिए यह मस्कोवाइट बनाता है।
\end{solution}

\begin{solution}{exo:b3:water-rock-interaction:10}
$\delta_w^f = -2 - 2 = -4$। बंद: $W/R = 0.5\times(-2 - 6)/(-13 + 4) =
0.44$। खुला: $0.5\ln[(-8)/(-9) + 1] = 0.5\ln1.889 = 0.32$। खुले तंत्र में प्रत्येक
पार्सल अपनी अंतिम भेंट वाली चट्टान के साथ साम्य में निकलता है और जितना संभव हो उतना
ऑक्सीजन-18 ले जाता है; बंद तंत्र में पूरा समूह $-4$ पर पहुँचता है और
प्रति किलोग्राम कम ले जाता है।
\end{solution}

\begin{solution}{exo:b3:water-rock-interaction:11}
खुला: $\log K = -5.59$, $m = (10^{-5.59}\times10^{-1.5}/4)^{1/3} =
\qty{2.73}{mmol/L}$, \qty{109}{mg/L}। बंद: $C_0 = 10^{-1.27 - 1.5} =
\qty{1.70e-3}{mol/L}$ और $K_{\mathrm{cal}}K_1/K_2 = 10^{-4.32}$; घन समीकरण
$4m^3 = 4.79\times10^{-5}(1.70\times10^{-3} - m)$ से $m =
\qty{1.45}{mmol/L}$, \qty{58}{mg/L}, लगभग आधा। शेष कार्बन डाइऑक्साइड
$0.25\times10^{-3}$ है, यानी $C_0$ का \qty{15}{\percent}, इसलिए $\log P_{\ce{CO2}}$
$-1.5$ से घटकर $-2.33$ हो जाता है।
\end{solution}

\begin{solution}{exo:b3:water-rock-interaction:12}
$1 - 10^{\mathrm{SI}}$: 0.99, 0.90 और 0.21। घात नियम देता है
$(10^4/6)^{-0.61} = 0.011$: दर लगभग 90 गुना कम। क्षेत्र का खनिज
पुराना है, उसकी अभिक्रियाशील सतहें समाप्त हो चुकी हैं, और वह साम्य के निकट धीमे छिद्र-जल
में डूबा है; ये दोनों कारक गुणा होते हैं, और क्षेत्रीय दर प्रयोगशाला दर से
कई कोटि नीचे गिर जाती है।
\end{solution}

\begin{solution}{pb:b3:water-rock-interaction:1}
\textbf{1.} \ce{Ca} 2.795, \ce{Mg} 0.165, \ce{Na} 0.305, \ce{K} 0.031,
\ce{HCO3} 5.572, \ce{Cl} 0.310, \ce{SO4} 0.146, \ce{NO3} \qty{0.145}{mmol/L}।

\textbf{2.} धनायन \qty{6.25}{meq/L}, ऋणायन \qty{6.32}{meq/L}: त्रुटि
\qty{-0.5}{\percent}, स्वीकार्य।

\textbf{3.} $I = \tfrac12(4\times2.795 + 4\times0.165 + 0.305 + 0.031 + 5.572
+ 0.310 + 4\times0.146 + 0.145)\times10^{-3} = 0.0094$।

\textbf{4.} $\sqrt I = 0.0969$: $\gamma_{\ce{Ca^2+}} = 0.683$,
$\gamma_{\ce{HCO3-}} = 0.904$, $\gamma_{\ce{CO3^2-}} = 0.676$।

\textbf{5.} हाँ: $I$ सीमा लगभग 0.1 का दसवाँ भाग है।

\textbf{6.} $a_{\ce{HCO3-}} = 5.572\times10^{-3}\times0.904 = 5.04\times
10^{-3}$ ($\log = -2.298$)।

\textbf{7.} $a_{\ce{CO2}} = 10^{-7.20}\times5.04\times10^{-3}/10^{-6.46} =
9.17\times10^{-4}$ ($\log = -3.038$)।

\textbf{8.} $\log P_{\ce{CO2}} = -3.038 + 1.27 = -1.77$, \qty{0.017}{atm}:
$10^{1.60} = 40$ गुना वायु के बराबर, जो भूजल परास के भीतर है। कार्बन डाइऑक्साइड
पुनर्भरण क्षेत्र की मृदा वायु से आती है, मूलों और अपघटन से।

\textbf{9.} $a_{\ce{CO3^2-}} = 10^{-10.49}\times5.04\times10^{-3}/10^{-7.20}
= 2.58\times10^{-6}$, $[\ce{CO3^2-}] = \qty{3.8e-6}{mol/L}$। कुल
\qty{6.49}{mmol/L}: \ce{CO2} \qty{14.1}{\percent}, \ce{HCO3-}
\qty{85.8}{\percent}, \ce{CO3^2-} \qty{0.06}{\percent}।

\textbf{10.} क्षारीयता $5.572 + 2\times0.0038 = \qty{5.580}{meq/L}$:
\qty{99.9}{\percent} बाइकार्बोनेट।

\textbf{11.} $a_{\ce{Ca^2+}} = 2.795\times10^{-3}\times0.683 = 1.909\times
10^{-3}$; $\log\mathrm{IAP} = \log(1.909\times10^{-3}\times2.58\times10^{-6})
= -8.307$; $\mathrm{SI} = -8.307 + 8.35 = +0.04$।

\textbf{12.} सूचकांक की अनिश्चितता $\pm0.1$ के भीतर है: जल
कैल्साइट के साथ साम्यावस्था में है। यह चट्टान में लंबे समय तक रहा है: एक धीमा, प्रसरित
सोता।

\textbf{13.} $m = [10^{-5.59}\times0.0171/(4\times0.683\times0.904^2)]^{1/3}
= \qty{2.70}{mmol/L}$, \qty{108}{mg/L}, मापे गए \qty{2.79}{mmol/L} के विरुद्ध:
\qty{3}{\percent} अधिक, जैसा कि $+0.04$ सूचकांक का अर्थ है।

\textbf{14.} $a_{\ce{CO2}} = 10^{-1.27 - 3.37} = 2.29\times10^{-5}$;
$a_{\ce{H+}} = 10^{-6.46}\times2.29\times10^{-5}/5.04\times10^{-3}$, pH 8.80;
$a_{\ce{CO3^2-}} = 1.03\times10^{-4}$; $\mathrm{SI} = \log(1.909\times
10^{-3}\times1.03\times10^{-4}) + 8.35 = +1.65$ (1.61, आकृति के पूर्ण
क्षारीयता संतुलन के साथ)।

\textbf{15.} पुनरावृत्ति करके $m = [10^{-5.59 - 3.37}/(4\gamma_{\ce{Ca}}
\gamma_{\ce{HCO3}}^2)]^{1/3}$, जहाँ $I = 3m$: $m = \qty{0.72}{mmol/L}$
(\qty{29}{mg/L}), $I = 0.0022$। कैल्साइट अवक्षेपित हुआ: $2.795 - 0.718 =
\qty{2.08}{mmol/L}$, लगभग \qty{210}{mg} प्रति लीटर।

\textbf{16.} $0.208\times2.0 = \qty{0.42}{g/s}$, \qty{13}{t} प्रति वर्ष,
$13.1/2.71 = \qty{4.8}{m^3}$; \qty{400}{m^2} पर, \qty{12}{mm} प्रति वर्ष,
\qty{0.03}{mm} प्रति दिन: अध्याय में उल्लिखित ठंडे सोतों के टूफ़ा से काफ़ी कम।

\textbf{17.} $X = \log(1.2\times10^{-3}/39.10) + 5.8 = 1.29$; $Y =
\log(9.0\times10^{-3}/60.08) = -3.82$।

\textbf{18.} $X < 1.55$ और $Y > -4.39$: कैओलिनाइट। क्वार्ट्ज़: $\mathrm{SI} =
-3.82 + 3.98 = +0.16$।

\textbf{19.} $\log a_{\ce{H4SiO4}} = -4.39$ से नीचे, \qty{2.4}{mg/L}
सिलिका।

\textbf{20.} $X = 1.55$ के लिए pH $= 1.55 + 4.51 = 6.06$ चाहिए।

\textbf{21.} अम्लीय $-10.06 - 0.500\times5.8 = -12.96$; उदासीन $-12.41$; क्षारीय
$-21.20 + 0.823\times5.8 = -16.43$। योग: $\log r = -12.30$, उदासीन पद
प्रमुख है।

\textbf{22.} आपूर्ति-सीमित: मोटी मृदा और मंद ढलान धीमे अपरदन का संकेत देते हैं;
अपक्षय अग्र तक पहुँचने वाला ताज़ा फ़ेल्सपार पूरी तरह निक्षालित हो जाता है, और
फ़्लक्स जलवायु के बजाय अपरदन दर का अनुसरण करता है।

\textbf{23.} $\exp[(74/0.008314)(1/283.15 - 1/298.15)] = 4.9$।

\textbf{24.} सोते के जल का कैल्साइट \omterm{def:b3:water-rock-interaction:si}{संतृप्ति सूचकांक} इसके निकास पर
है $\mathrm{SI} = \textbf{+0.04}$।
\end{solution}
